\section{Discrete resonator lattices and finite arrays}\label{sec:discrete_lattice}

In the previous Section~\ref{sec:dispersion_relation}, we have showed that in the local effective-medium model the fluid added mass dilutes the resonator mass ratio, equation~\eqref{eq:lat_mu}. Therefore, the homogenised LRH structure has no complete stop band.  With this spatial averaging, the discrete geometry of the resonator array is not considered and, therefore, all host Fourier components generated by the attachment pattern are neglected. In a practical setting, and also in certain wave regimes, this omission might be significant, because the fluid added mass $M_{f,\rho}$ in equation~\eqref{eq:plane_wave_added_mass} diverges as $|\mathbf{k}|\to0$. That is, the components discarded by homogenization therefore control the band structure at long wavelength (Section~\ref{subsubsec:lat_long_wave}). In this section we keep the resonators discrete and represent the bare hydroelastic host exactly by its Green function. The aim is to determine what sets band gaps, scattering and absorption once this averaging is removed.

Here we use classical tools in structural dynamics literature, i.e. Bloch lattice sums and multiple scattering for point scatterers and resonators on plates \cite{evans2007penetration,movchan2007bloch,xiao2012flexural,torrent2013elastic}, the Foldy multiple-scattering system \cite{foldy1945multiple,martin2006multiple}, and wave-energy absorption bounds \cite{budal1975resonant,evans1976theory,falnes1980radiation}. Graded resonator arrays on elastic plates loaded by a compressible fluid were studied in \cite{skelton2018multiphysics}; here the fluid has a free surface and loads the plate through gravity waves. For the sake of completeness, we restate them for the floating structure Green function, where the fluid enters through the host symbol $\mathcal{D}_0$, through its wavenumber-dependent added mass and, for finite structures, through radiation into open water. The main novel results discussed in this section regarding the unique properties of LRH structures are threefold. First, a complete stop band of a floating lattice, when it exists, is set by the discrete geometry through the lattice near field, and not by the effective mass (Section~\ref{subsubsec:lat_long_wave}). Second, the radiative loss rate of a resonator is defined by the hydroelastic host, $\gamma_r=k_0/(2\mathcal{D}_0'(k_0))$ (Section~\ref{subsubsec:energy_flux_single}), determining the critical damping for maximum absorption. {\color{red}Third, for a finite structure, the cut-off of the added-mass divergence makes the attainable hybridisation depend on the size of the structure (Section~\ref{subsubsec:lat_long_wave}) [pending verification against the finite-element model].}

In what follows, we consider two configurations: an infinite floating structure, denoted by \emph{configuration~A}, for which Bloch theory applies; and a finite structure in open water, denoted by \emph{configuration~B}, which provides a direct link to the finite element formulation developed in Section~\ref{sec:numerical_formulation}.

\subsection{Single resonator on infinite structure}\label{subsec:single_resonator}

Let us consider the host symbol $\mathcal{D}_0(\mathbf{k},\omega)$, leading to the dispersion relation, of the bare host plate, i.e. without effect of the resonators. Starting from equation~\eqref{eq:dispersion_relation}, with $m_{\mathrm{eff},\rho}\to m_\rho$, multiplied by $\rho$, and using the dimensional plate mass $m=\rho m_\rho$, tension $T=\rho T_\rho$ and bending stiffness $D=\rho C_\rho$, we have
\begin{equation}
  \mathcal{D}_0(\mathbf{k},\omega)
  = \rho g + (1-\mathrm{i}\omega\tau)\,Tk^2
  + (1-\mathrm{i}\omega\tau_b)\,Dk^4
  - \omega^2\Big[m+\frac{\rho}{k\tanh kH}\Big],
  \qquad k=|\mathbf{k}| .
  \label{eq:host_symbol}
\end{equation}
For an anisotropic plate, the bending term is replaced by the corresponding tensor, as discussed in Section~\ref{sec:dispersion_relation}. The last term is the fluid added mass \eqref{eq:plane_wave_added_mass}. It does not depend on $\omega$ but diverges as $k\to0$, which is the key property discussed in the following sections.

From equation~\eqref{eq:resonator_force}, a resonator of mass $m_r$, spring constant $k_r$ and damping constant $c_r$ attached at $\mathbf{x}_j$ exerts the force $F_j=\alpha_j\widehat{\eta}_j$ towards the plate, with
$\widehat{\eta}_j=\widehat{\eta}(\mathbf{x}_j)$ the bare plate deflection at the attachment point, and attachment strength $\alpha_j=\omega^2m_{\mathrm{eff},j}$.  The equivalent compliance form leads to
\begin{equation}
  \label{eq:alpha_j}
  \frac{1}{\alpha_j}=\frac{1}{\omega^2 m_r}-\frac{1}{k_r^\omega},
  \qquad k_r^\omega=k_r-\mathrm{i}\omega c_r ,
\end{equation}
where the compliances of the mass and of the spring--dashpot are added in series. Note that only the damping contributes to the imaginary part, i.e.
$\gamma_d\equiv-\operatorname{Im}(1/\alpha_j)=\omega c_r/|k_r^\omega|^2\ge0$.

Let us denote by $\mathcal{L}_0$ the bare host operator, i.e. the force per
unit area required to hold the fluid-loaded (floating) structure in the deflection $\widehat{\eta}$. Because of the added-mass term, $\mathcal{L}_0$ is a non-local operator. Then, the structure is forced by an applied pressure $\widehat{f}$ and by the resonator forces $F_j=\alpha_j\widehat{\eta}_j$, each concentrated at its attachment point, leading to an implicit equation for the deflection
\begin{equation}
  \label{eq:lattice_coupled}
  \mathcal{L}_0\widehat{\eta}(\mathbf{x})
  = \widehat{f}(\mathbf{x})
  + \sum_j \alpha_j(\omega)\,\widehat{\eta}_j\,
    \delta(\mathbf{x}-\mathbf{x}_j).
\end{equation}

Equation~\eqref{eq:lattice_coupled} can be inverted using the Green function of the structure, $G$, i.e. the deflection of the floating structure under a unit point force applied at the origin, $\mathcal{L}_0G=\delta(\mathbf{x})$. Assuming an infinite structure, \emph{configuration A}, the floating structure is translation invariant, therefore, a unit force at $\mathbf{x}_j$ produces the deflection $G(\mathbf{x}-\mathbf{x}_j)$. By superposition, the solution of \eqref{eq:lattice_coupled} reads
\begin{equation}
  \label{eq:green_solution}
  \widehat{\eta}(\mathbf{x})
  = \widehat{\eta}_f(\mathbf{x})
  + \sum_j \alpha_j(\omega)\,\widehat{\eta}_j\,G(\mathbf{x}-\mathbf{x}_j),
\end{equation}
where $\widehat{\eta}_f$ is the response of the bare floating structure to the applied pressure $\widehat{f}$, e.g. from an incident wave. Note that equation~\eqref{eq:green_solution} is still implicit in the attachment deflections $\widehat{\eta}_j$, but evaluating it at the attachment points reduces the problem to a set of algebraic equations, which is the starting point of the single-resonator, lattice and finite-array analyses below.

Since $\mathcal{L}_0$ acts in Fourier space as a multiplication by $\mathcal{D}_0(\mathbf{k},\omega)$, the Fourier transform of the Green function is simply $1/\mathcal{D}_0(\mathbf{k},\omega)$. For an isotropic plate, $\mathcal{D}_0$ depends only on $k=|\mathbf{k}|$, $G$ depends only on $r=|\mathbf{x}|$, and the inverse Fourier transform reduces to a Hankel transform. This transform can be evaluated by residues, leading to the classical representation of the floating-plate Green function on water of finite depth, e.g. \cite{evans2003wave}
\begin{equation}
  \label{eq:green_function}
  G(r)=\frac{1}{2\pi}\int_0^\infty
       \frac{J_0(kr)\,k}{\mathcal{D}_0(k,\omega)}\,\mathrm{d}k
      =\frac{\mathrm{i}}{2}\sum_{n}
       \frac{k_n}{\mathcal{D}_0'(k_n)}\,H_0^{(1)}(k_nr),
  \qquad \mathcal{D}_0'\equiv\partial_k\mathcal{D}_0.
\end{equation}
The sum runs over the roots of the bare dispersion relation, $\mathcal{D}_0(k_n,\omega)=0$, with $\operatorname{Im}k_n\ge0$, and each root contributes one type of wave to the response \citep{fox1994oblique}: the real root $k_0$ gives the outgoing flexural--gravity wave that carries energy to infinity; a complex pair, associated with plate flexure, gives an oscillation that decays away from the load; and an infinite sequence of imaginary roots, $k_n\simeq\mathrm{i}n\pi/H$, gives evanescent fluid modes confined within a distance of order $H$ from the load.

In the non-dissipative case, only the propagating wave contributes to the imaginary part of $G$, since the complex pair and the imaginary roots contribute real terms. Hence
\begin{equation}
  \label{eq:ImG}
  \operatorname{Im}G(r)=\frac{k_0}{2\mathcal{D}_0'(k_0)}\,J_0(k_0r),
  \qquad
  \gamma_r\equiv\operatorname{Im}G(0)=\frac{k_0}{2\mathcal{D}_0'(k_0)}>0.
\end{equation}
The quantity $\gamma_r$ is the radiative counterpart of $\gamma_d$, it measures the power radiated as waves by a point force at a given point of application. Finally, since $\mathcal{D}_0\sim Dk^4$ for large $k$, the integral in \eqref{eq:green_function} converges at $r=0$ and $G(0)$ is finite for any $D>0$. Therefore, a point attachment is well posed whenever its footprint radius $r_c$ is small compared with both the lattice spacing and the flexural wavelength. 
%A finite circular footprint can be included by replacing $1/\mathcal{D}_0$ with $\mathcal{F}^2(kr_c)/\mathcal{D}_0$, where $\mathcal{F}(s)=2J_1(s)/s$, in all wavenumber integrals and sums. This regularisation is only needed in the tension-dominated limit $D\to0$.

\subsubsection{Energy flux of a single resonator}
\label{subsubsec:energy_flux_single}

To quantify the power scattered and absorbed by the resonators, and later to define reflection and transmission coefficients, we need the energy flux carried by a plane flexural--gravity wave. Rather than computing it from a direct energy balance perspective, here we derive the flux from the Green function by a power balance on a single point force. Let us consider a point force $F$ applied at the origin of the bare structure. From equation~\eqref{eq:green_solution}, with no resonators and $\widehat{\eta}_f=0$, the deflection at the point of application is $\widehat{\eta}(0)=G(0)F$, and the corresponding velocity is $-\mathrm{i}\omega\widehat{\eta}(0)$. Using standard time averaging principles of time-harmonic quantities, we obtain the power injected by the force as
\begin{equation}
  \label{eq:point_power}
  P=\tfrac12\operatorname{Re}\!\big[F\,\overline{(-\mathrm{i}\omega G(0)F)}\big]=\tfrac12\,\omega|F|^2\operatorname{Im}G(0)=\tfrac12\,\omega|F|^2\gamma_r .
\end{equation}
The real part of $G(0)$ gives a response in phase with the displacement, which stores and returns energy over a cycle, and only the imaginary part, in phase with the velocity, does net work. This justifies the interpretation of $\gamma_r$ given above. Since the bare structure is non-dissipative, this power cannot be lost in the structure or in the fluid and must leave the load as outgoing waves.

Matching this power to the flux of the outgoing far field of the Green function~\eqref{eq:green_function} (Appendix~\ref{app:flux_derivation}) gives the time-averaged energy flux per unit crest length of a plane flexural--gravity wave of amplitude $A$,
\begin{equation}
  \label{eq:energy_flux}
  I=\frac{\omega}{4}\,\mathcal{D}_0'(k_0,\omega)\,|A|^2 .
\end{equation}

\begin{remark}
  For a membrane without fluid, $\mathcal{D}_0=Tk^2-m\omega^2$ and $\mathcal{D}_0'(k_0)=2Tk_0$, so that equation~\eqref{eq:energy_flux} recovers the classical flux $\tfrac12T\omega k_0|A|^2$. In the hydroelastic case, the symbol $\mathcal{D}_0$ is obtained after eliminating the fluid, and the kinetic energy of the fluid is contained in it through the added-mass term. Therefore, equation~\eqref{eq:energy_flux} accounts consistently for the shear and moment work of the structure and for the pressure work of the fluid, without computing them separately.
\end{remark}

\begin{remark}
  Equation~\eqref{eq:energy_flux} confirms the sign of $\gamma_r$ in equation~\eqref{eq:ImG}. In the non-dissipative case, $\mathcal{D}_0(k,\omega)\to-\infty$ as $k\to0$, because the fluid added mass diverges, and $\mathcal{D}_0(k,\omega)\to+\infty$ as $k\to\infty$, because of the bending term. Since $k_0$ is the only positive real root of $\mathcal{D}_0$ \citep{fox1994oblique}, $\mathcal{D}_0$ changes sign from negative to positive at $k_0$, so that $\mathcal{D}_0'(k_0)>0$. Hence the flux \eqref{eq:energy_flux} is positive, the propagating term in \eqref{eq:green_function} indeed carries energy away from the load, and $\gamma_r>0$.
\end{remark}

With the Green function and the energy flux at hand, we can now analyse the response of a single resonator attached at the origin of the infinite structure and excited by the incident plane wave of amplitude $A$. The incident wave induces $\widehat{\eta}_f$ in equation~\eqref{eq:green_solution}, and the resonator acts on the structure as a point force $F_1=\alpha\widehat{\eta}_1$. Hence, the total deflection is
\begin{equation}
  \label{eq:single_field}
  \widehat{\eta}(\mathbf{x})=A\mathrm{e}^{\mathrm{i}k_0\hat{\mathbf{e}}\cdot\mathbf{x}}+G(|\mathbf{x}|)\,F_1 .
\end{equation}
Evaluating equation~\eqref{eq:single_field} at the attachment point gives $\widehat{\eta}_1=A+\alpha G(0)\widehat{\eta}_1$. That is, the resonator is not driven by the incident wave alone, but by the incident wave plus its own field at the attachment point, $G(0)F_1$. Solving for $\widehat{\eta}_1$, the resonator force reads
\begin{equation}
  \label{eq:single_force}
  F_1=\alpha\widehat{\eta}_1=\frac{A}{1/\alpha-G(0)},
  \qquad
  \frac{1}{\alpha}-G(0)=x-\mathrm{i}(\gamma_d+\gamma_r),
\end{equation}
where we have used $\operatorname{Im}(1/\alpha)=-\gamma_d$ from equation~\eqref{eq:alpha_j} and $\operatorname{Im}G(0)=\gamma_r$ from equation~\eqref{eq:ImG}. The denominator has the structure of a damped oscillator. Its real part, $x=\operatorname{Re}(1/\alpha)-\operatorname{Re}G(0)$, is a detuning that compares the compliance of the resonator with the point compliance of the floating structure. Its imaginary part collects the two loss channels of the resonator: the dissipation in the dashpot, $\gamma_d$, and the radiation of waves into the structure, $\gamma_r$.

The scattered and absorbed powers follow from equation~\eqref{eq:single_force}. The scattered field, $G(|\mathbf{x}|)F_1$, is exactly the field of a point force in the bare structure, so the power it carries to infinity is given by equation~\eqref{eq:point_power} with $F=F_1$, i.e. $P_{\mathrm{sc}}=\tfrac12\omega|F_1|^2\gamma_r$. The absorbed power is the work done by the structure on the resonator. By action and reaction, the structure exerts the force $-F_1$ on the resonator, whose attachment point moves with velocity $-\mathrm{i}\omega\widehat{\eta}_1=-\mathrm{i}\omega F_1/\alpha$. Hence, $P_{\mathrm{abs}}=\tfrac12\operatorname{Re}[-F_1\overline{(-\mathrm{i}\omega F_1/\alpha)}]=\tfrac12\omega|F_1|^2\gamma_d$, which is the power dissipated in the dashpot. Substituting $F_1$ from equation~\eqref{eq:single_force}, we obtain
\begin{equation}
  \label{eq:single_powers}
  P_{\mathrm{sc}}=\frac{\omega|A|^2}{2}\,\frac{\gamma_r}{x^2+(\gamma_d+\gamma_r)^2},
  \qquad
  P_{\mathrm{abs}}=\frac{\omega|A|^2}{2}\,\frac{\gamma_d}{x^2+(\gamma_d+\gamma_r)^2}.
\end{equation}
Dividing the two powers by the incident flux per unit crest length we obtain the scattering and absorption cross-sections, also known as capture width. Using $1/\mathcal{D}_0'(k_0)=2\gamma_r/k_0$ from equation~\eqref{eq:ImG}, the scattering and absorption cross-sections read
\begin{equation}
  \label{eq:single_sigma}
  \sigma_{\mathrm{sc}}=\frac{4}{k_0}\,\frac{\gamma_r^2}{x^2+(\gamma_d+\gamma_r)^2},
  \qquad
  \sigma_{\mathrm{abs}}=\frac{4}{k_0}\,\frac{\gamma_r\gamma_d}{x^2+(\gamma_d+\gamma_r)^2}.
\end{equation}
Two results follow from equation~\eqref{eq:single_sigma}, a renormalisation of the resonance frequency and an upper bound on the absorbed power. These are reflected in the following remarks.

\begin{remark}[Renormalised resonance] In the undamped case, equation~\eqref{eq:alpha_j} gives $1/\alpha=1/(\omega^2m_r)-1/k_r$, and this condition becomes
\begin{equation}
  \label{eq:renorm_single}
  \widetilde{\omega}_r^{\,2}=\frac{\widetilde{k}_r}{m_r},
  \qquad
  \frac{1}{\widetilde{k}_r}=\frac{1}{k_r}+\operatorname{Re}G(0;\omega).
\end{equation}
That is, the resonator spring acts in series with the point compliance of the floating structure, $\operatorname{Re}G(0)$, which itself depends on frequency, so that equation~\eqref{eq:renorm_single} is an implicit equation for $\widetilde{\omega}_r$. The natural frequency of the resonator in isolation $\omega_r=\sqrt{k_r/m_r}$, used in the effective-medium model of Section~\ref{sec:dispersion_relation}, is recovered only when the spring is much softer than the structure at the attachment point, $k_r\operatorname{Re}G(0)\ll1$.
\end{remark}

\begin{remark}[Absorption bound] The absorption cross-section, $\sigma_{\mathrm{abs}}$, vanishes both for $\gamma_d=0$, where nothing is dissipated, and for $\gamma_d\gg\gamma_r$, where dissipation dominates the denominator and damps out the resonator response. Its maximum is attained at critical coupling, $\gamma_d=\gamma_r$, i.e. when the dissipative and radiative loss rates are equal, leading to
\begin{equation}
  \label{eq:absorption_bound}
  \sigma_{\mathrm{abs}}\le\sigma_{\mathrm{abs}}^{\max}=\frac{1}{k_0}=\frac{\lambda}{2\pi}.
\end{equation}
At critical coupling the resonator absorbs as much power as it scatters, $P_{\mathrm{abs}}=P_{\mathrm{sc}}$. The bound \eqref{eq:absorption_bound} depends only on the wavelength, whereas the damping that attains it depends, through $\mathcal{D}_0'(k_0)$, on the structure, the fluid depth and the pre-tension. Using $\gamma_d=\omega c_r/|k_r^\omega|^2$, the critical-coupling condition reads
\begin{equation}
  \label{eq:critical_coupling}
  \frac{\omega c_r}{k_r^2+\omega^2c_r^2}=\frac{k_0}{2\mathcal{D}_0'(k_0)} .
\end{equation}
Since the left-hand side is bounded by $1/(2k_r)$, its value at $\omega c_r=k_r$, equation~\eqref{eq:critical_coupling} has a real solution for $c_r$ only if $k_r\gamma_r\le1/2$. Otherwise, the radiative width exceeds any attainable dissipative width and the bound \eqref{eq:absorption_bound} cannot be reached. Equation~\eqref{eq:absorption_bound} is the hydroelastic counterpart of the capture-width limit $\lambda/(2\pi)$ for an axisymmetric heaving point absorber in open water \citep{budal1975resonant,evans1976theory}. In the present setting, however, the waves are flexural--gravity waves of the floating structure, and the resonator radiates only through the structure.
\end{remark}
\paragraph{Radiation damping in the vertical section.} In the vertical section, $\gamma_r=1/\mathcal{D}_0'(k_0)$ (equation~\eqref{eq:lat_green_1d}), and the critical-coupling condition~\eqref{eq:critical_coupling} evaluated at resonance, where $\omega c_r=2\zeta k_r$, gives the optimal damping ratio
\begin{equation}\label{eq:zeta_opt}
    \frac{2\zeta_{\rm opt}}{1+4\zeta_{\rm opt}^2}=k_r\gamma_r=\frac{m_{\rho,r}\,\omega_r^2}{\mathcal{D}_0'(k_0)},
    \qquad
    \zeta_{\rm opt}\simeq\frac{m_{\rho,r}\,\omega_r^2}{2\,\mathcal{D}_0'(k_0)},
\end{equation}
with $m_{\rho,r}$ per unit width and $\mathcal{D}_0'$ the density-normalised derivative of equation~\eqref{eq:weak_dk}. Since $|\partial_kM_{f,\rho}(k,H)|\ge1/k^2$ at any depth, $\mathcal{D}_0'(k_0)\ge\omega^2/k_0^2$, and
\begin{equation}\label{eq:zeta_bound}
    \zeta_{\rm opt}\lesssim\frac{m_{\rho,r}\,k_0^2}{2}.
\end{equation}
The optimal damping, and with it the absorption bandwidth at critical coupling, $\Delta\omega/\omega_r\simeq4\zeta_{\rm opt}$, is therefore set by the resonator mass per unit width compared with the fluid area $1/k_0^2$. For the parameters of Table~\ref{tab:dispersion_test_parameters} and $\omega_r=10~\si{\radian\per\second}$, $\zeta_{\rm opt}=0.0093$ for $T_\rho=0$ and $0.0064$ for $T_\rho>0$, with the fluid contributing $25\%$ and $34\%$ of $\mathcal{D}_0'(k_0)$, respectively. At $\omega_r=1~\si{\radian\per\second}$, the fluid contributes more than $98\%$ of $\mathcal{D}_0'(k_0)$, and $\zeta_{\rm opt}=4.3\times10^{-5}$ for both pre-tensions (Figure~\ref{fig:lattice_design}(b)). A single resonator of fixed mass therefore becomes an extremely narrow-band absorber at low frequency, a direct consequence of the fluid inertia, and the same dilution identified in Section~\ref{subsubsec:fluid_dilution} for the hybridisation. {\color{red}The finite element model reproduces this collapse: the damping ratio that maximises the absorption falls from $9.26\times10^{-3}$ at $\omega_r=10~\si{\radian\per\second}$ to $4.35\times10^{-5}$ at $\omega_r=1~\si{\radian\per\second}$ ($T_\rho=0$), and the bound $\mathcal{A}=\tfrac12$ is reached at critical coupling at every frequency (Section~\ref{subsubsec:num_L1_3b}).}

\subsection{Infinite array of resonators on infinite structure}
\label{subsec:infinite_array}

We now consider identical resonators, characterized by $\alpha_j=\alpha$, attached at the nodes, denoted by $\mathbf{R}$, of a Bravais lattice $\Lambda$ with primitive cell area $A_c$. We denote by $\mathbf{K}$ the vectors of the reciprocal lattice $\Lambda^*$. Free waves of the lattice are solutions of equation~\eqref{eq:green_solution} without external forcing, $\widehat{\eta}_f=0$, with Bloch form, $\widehat{\eta}(\mathbf{x}+\mathbf{R})=\mathrm{e}^{\mathrm{i}\mathbf{q}\cdot\mathbf{R}}\widehat{\eta}(\mathbf{x})$, and Bloch wavevector $\mathbf{q}$. The attachment deflections are then $\widehat{\eta}_{\mathbf{R}}=\mathrm{e}^{\mathrm{i}\mathbf{q}\cdot\mathbf{R}}\widehat{\eta}_{\mathbf{0}}$, and evaluating equation~\eqref{eq:green_solution} at the origin gives $\widehat{\eta}_{\mathbf{0}}=\alpha\,\widehat{\eta}_{\mathbf{0}}\sum_{\mathbf{R}}G(\mathbf{R})\mathrm{e}^{\mathrm{i}\mathbf{q}\cdot\mathbf{R}}$. A non-trivial solution exists if and only if
\begin{equation}
  \label{eq:bloch_condition}
  \frac{1}{\alpha(\omega)}=\mathcal{S}(\mathbf{q},\omega),
  \qquad
  \mathcal{S}(\mathbf{q},\omega)\equiv\sum_{\mathbf{R}\in\Lambda}G(\mathbf{R})\,\mathrm{e}^{\mathrm{i}\mathbf{q}\cdot\mathbf{R}}
  =\frac{1}{A_c}\sum_{\mathbf{K}\in\Lambda^*}\frac{1}{\mathcal{D}_0(\mathbf{q}+\mathbf{K},\omega)} ,
\end{equation}
where the second equality follows from the Poisson summation formula and from $1/\mathcal{D}_0$ being the Fourier transform of $G$. Equation~\eqref{eq:bloch_condition} is the exact dispersion relation of the lattice that generalises the single-resonator condition of equation~\eqref{eq:single_force}. Here, the self-interaction $G(0)$ is replaced by the lattice sum $\mathcal{S}$, containing the self term $\mathbf{R}=\mathbf{0}$ plus the field produced at the origin by all the other resonators. Since $G(r)$ decays only like $r^{-1/2}$, the real-space sum is not absolutely convergent, and the reciprocal-space form is used as the definition of $\mathcal{S}$. Its terms decay like $1/(D|\mathbf{K}|^4)$, thus, the sum converges for any $D>0$.

In the non-dissipative case, the lattice does not radiate. The terms of the reciprocal sum in equation~\eqref{eq:bloch_condition} are real, except on the empty-lattice circles $|\mathbf{q}+\mathbf{K}|=k_0$, where $\mathcal{D}_0$ vanishes. Equivalently, by equation~\eqref{eq:ImG}, $\operatorname{Im}\mathcal{S}=\gamma_r\sum_{\mathbf{R}}J_0(k_0|\mathbf{R}|)\mathrm{e}^{\mathrm{i}\mathbf{q}\cdot\mathbf{R}}$, which by Poisson summation is a sum of delta functions supported on these circles. Away from them, $\mathcal{S}$ and $1/\alpha$ are both real, and equation~\eqref{eq:bloch_condition} is a real dispersion relation. In contrast with the single resonator, in a periodic array, the waves scattered by the individual resonators interfere coherently to form the Bloch wave, and no energy is lost to infinity.

The relation between equation~\eqref{eq:bloch_condition} and the homogenised law of Section~\ref{sec:dispersion_relation} becomes explicit by separating the $\mathbf{K}=\mathbf{0}$ term of the reciprocal sum. Multiplying equation~\eqref{eq:bloch_condition} by $A_c$ and rearranging gives the exact identity
\begin{equation}
  \label{eq:exact_em}
  \mathcal{D}_0(\mathbf{q},\omega)=\frac{\widetilde{\alpha}(\mathbf{q},\omega)}{A_c},
  \qquad
  \frac{1}{\widetilde{\alpha}}=\frac{1}{\omega^2m_r}-\frac{1}{k_r^\omega}-\frac{\mathcal{S}'(\mathbf{q},\omega)}{A_c},
  \qquad
  \mathcal{S}'(\mathbf{q},\omega)=\sum_{\mathbf{K}\neq\mathbf{0}}\frac{1}{\mathcal{D}_0(\mathbf{q}+\mathbf{K},\omega)} .
\end{equation}
For $\widetilde{\alpha}=\alpha$, equation~\eqref{eq:exact_em} is precisely the homogenised dispersion relation \eqref{eq:dispersion_relation} for identical resonators with number density $n=1/A_c$, written in dimensional form, $\mathcal{D}_0=n\alpha$. The $\mathbf{K}=\mathbf{0}$ term represents the mean force exerted by the resonators, which is all that homogenisation retains, whereas $\mathcal{S}'$ collects the components of the field that vary on the scale of the lattice and that homogenisation discards. When the lattice is subwavelength, $|\mathbf{q}+\mathbf{K}|>k_0$ for all $\mathbf{K}\neq\mathbf{0}$, all the terms of $\mathcal{S}'$ are positive and the near field softens the effective spring, lowering the resonance.

\begin{remark}[Exact effective-medium form]
Homogenisation is therefore exact in the limit $\mathcal{S}'\to0$, and its accuracy is measured by the ratio between the near-field compliance and the spring compliance. For a square lattice of spacing $a$, the reciprocal vectors have magnitude $|\mathbf{K}|\ge2\pi/a$. If bending dominates the symbol at this scale, $1/\mathcal{D}_0(\mathbf{K})\simeq1/(D|\mathbf{K}|^4)$, and this ratio, evaluated at $\mathbf{q}=\mathbf{0}$, reads
\begin{equation}
  \label{eq:PiL}
  \Pi_L=\frac{k_r\,\mathcal{S}'(\mathbf{0},\omega_r)}{A_c}\simeq\frac{k_r\,a^2}{16\pi^4D}\,\mathcal{Z}(2)\approx3.87\times10^{-3}\,\frac{k_ra^2}{D},
\end{equation}
where $\mathcal{Z}(2)=\sum_{\mathbf{n}\in\mathbb{Z}^2\setminus\{\mathbf{0}\}}|\mathbf{n}|^{-4}=4\zeta(2)\beta(2)\approx6.027$ is the Epstein zeta function of the square lattice, with $\zeta$ and $\beta$ the Riemann zeta and Dirichlet beta functions. The lattice-compliance number $\Pi_L$ compares the stiffness of the resonator spring with the bending stiffness of the structure over one lattice cell. To leading order, the resonance frequency entering the homogenised law is shifted to $\widetilde{\omega}_r\simeq\omega_r(1+\Pi_L)^{-1/2}$. Importantly, $\Pi_L$ does not depend on the wavelength. The scale-separation condition \eqref{eq:homogenization_scale_separation}, which for a lattice reads $|\mathbf{q}|a\ll\pi$, is therefore necessary but not sufficient: the effective-medium model is asymptotically exact only if, in addition, $\Pi_L\ll1$. Stiff resonators on a flexible structure can violate this condition even for a deeply subwavelength lattice.
\end{remark}

{\color{red}
\paragraph{Stop bands.}
For real $\omega$ and lossless constituents, a complete stop band is a frequency interval in which equation~\eqref{eq:bloch_condition} has no root for any real $\mathbf{q}$ in the first Brillouin zone. Since band extrema need not lie on the boundary of the irreducible Brillouin zone, the whole zone must be examined unless this is shown otherwise. With damping, we compute complex Bloch wavevectors $\mathbf{q}(\omega)$ at real $\omega$, whose imaginary part gives the spatial attenuation rate relevant to forced finite arrays (Section~\ref{subsec:finite_array}). Complex frequencies $\omega(\mathbf{q})$ at real $\mathbf{q}$ describe free temporal decay instead, and are not used to infer transmission.
}

\subsubsection{Long-wave limit and complete stop bands}\label{subsubsec:lat_long_wave}

In configuration~A, i.e. infinite floating structure, the fluid added mass \eqref{eq:plane_wave_added_mass} diverges at long wavelength, $M_{f,\rho}(k,H)\sim1/(k^2H)$ at finite depth and $M_{f,\rho}\sim1/k$ in deep water. Hence, the host symbol \eqref{eq:host_symbol} satisfies $\mathcal{D}_0(\mathbf{q},\omega)\to-\infty$ as $\mathbf{q}\to\mathbf{0}$, for every $\omega>0$ and any depth. This is the spectral counterpart of incompressibility, that is, a free surface fully covered by the structure cannot heave uniformly without changing the volume of the fluid below it. Both the homogenised law, $\mathcal{D}_0=\alpha/A_c$, and the exact law \eqref{eq:exact_em}, $\mathcal{D}_0=\widetilde{\alpha}/A_c$, must accommodate this divergence. However, this leads to two distinct results, which we analyse in what follows.

\paragraph{Homogenised model.} With $\mathcal{D}_0=\alpha/A_c$, the divergence can only be accommodated by $1/\alpha\to0$, i.e. $\omega\to\omega_r$ in the undamped case. This is the limit $s_+(k)\to\omega_r^2$ of Section~\ref{subsubsec:fluid_dilution}: the mass ratio \eqref{eq:lat_mu}, with $n_0=1/A_c$, vanishes at long wavelength, and the homogenised model has no complete stop band for any choice of parameters.

\paragraph{Discrete lattice.} In the exact law \eqref{eq:exact_em}, the same divergence requires $1/\widetilde{\alpha}(\mathbf{q},\omega)\to0$ as $\mathbf{q}\to\mathbf{0}$. Therefore, the band resulting from the centre $\Gamma$ of the Brillouin zone above the lower band starts at the frequency $\omega_\Gamma$ given, in the undamped case, by
\begin{equation}\label{eq:lat_gamma}
    \frac{1}{\omega_\Gamma^2m_r}=\frac{1}{k_r}+\frac{\mathcal{S}'(\mathbf{0},\omega_\Gamma)}{A_c},
    \qquad
    \omega_\Gamma\simeq\frac{\omega_r}{\sqrt{1+\Pi_L}} .
\end{equation}
That is, the upper band edge is the resonance renormalised by the lattice near field, which homogenisation discards. {\color{red}The estimate on the right of equation~\eqref{eq:lat_gamma} holds for $2\pi/a>k_c$, where the folded host branches lie above $\omega_r$ and $\Pi_L>0$; for $2\pi/a<k_c$, $\Pi_L<0$ and it fails. Its error is formally $O(\Pi_L^2)$. When bending dominates at the lattice scale it behaves as $\Pi_L^3$, because the coefficient of $\Pi_L^2$, which comes from evaluating $\mathcal{S}'$ at $\omega_r$ instead of $\omega_\Gamma$, is itself proportional to $\Pi_L$ (Section~\ref{subsubsec:num_L1_7}).} On the other side, the lower band terminates at the first Brillouin zone boundary, $\mathbf{q}_X$, at the frequency $\omega_X$ given by the root of $1/\alpha(\omega)=\mathcal{S}(\mathbf{q}_X,\omega)$, which lies below $\omega_r$. Therefore, there exists a complete stop band $[\omega_X,\omega_\Gamma]$ whenever $\omega_X<\omega_\Gamma$, i.e. when the depression of the lower band at the zone boundary exceeds the near-field shift of the resonance at $\Gamma$. Both edges are set by the discrete geometry, through $\mathcal{S}'(\mathbf{0},\omega)$ and $\mathcal{S}(\mathbf{q}_X,\omega)$, and the existence of the stop band cannot be inferred from the sign of the effective inertia of the homogenised model. {\color{red}For the parameters of Table~\ref{tab:dispersion_test_parameters}, \cite{liu2025metaplate} computed the exact bands numerically and reported a gap of about $9.5$--$10~\si{\radian\per\second}$, consistent with the values below, and interpreted it with the local branches restricted to the first Brillouin zone, i.e. with upper edge at $\omega_r$ and lower edge at $\omega_-(\pi/a)$. The lattice sums show instead that, without this truncation, the local model has no complete stop band, and that both edges are set by the discrete geometry: the upper edge by the near-field renormalisation~\eqref{eq:lat_gamma}, below $\omega_r$, and the lower edge by the zone-folded component $\mathbf{K}=-2\pi/a$.}

Figure~\ref{fig:lat_bands} compares the exact bands, computed from equation~\eqref{eq:bloch_condition}, with the local branches \eqref{eq:omega_k_x_closed_form} for a line array with the parameters of Table~\ref{tab:dispersion_test_parameters} and spacing $a=1~\si{\meter}$. For a line array, the formulation of Section~\ref{subsec:infinite_array} applies with $A_c=a$ and $\mathbf{K}=2\pi n/a$, $n\in\mathbb{Z}$, and the bending-only estimate of the lattice-compliance number becomes $\Pi_L\simeq k_ra^3/(720D)$, since $\sum_{n\neq0}n^{-4}=\pi^4/45$. Away from the Brillouin zone boundary, $\Pi_L=2.86\times10^{-3}$ for $T_\rho=0$, close to the bending-only estimate $2.83\times10^{-3}$, and $\Pi_L=1.94\times10^{-3}$ for $T_\rho=9.8~\si{\meter\cubed\per\second\squared}$, for which the pre-tension contributes to the host symbol at the lattice scale and stiffens the near field.

The exact lattice has a complete stop band $[\omega_X,\omega_\Gamma]=[9.646,\,9.986]~\si{\radian\per\second}$ ($\Delta\omega/\omega_r=0.034$) for $T_\rho=0$, and $[9.916,\,9.990]~\si{\radian\per\second}$ ($\Delta\omega/\omega_r=0.0074$) for $T_\rho=9.8~\si{\meter\cubed\per\second\squared}$. The anticrossing separations of Figure~\ref{fig:subsubsec_non_dissipative_homogeneous_isotropic_dispersion}, $0.163$ and $0.138$, overestimate these widths by factors of about $5$ and $19$, respectively. Truncating the local branches at $qa=\pi$ would give $[9.813,\,10]$ and $[9.959,\,10]~\si{\radian\per\second}$, misplacing both edges: the upper edge because the near-field shift of the resonance is missing, and the lower edge because, at the zone boundary $q=\pi/a$, the components $K=0$ and $K=-2\pi/a$ have the same wavenumber magnitude, $|q+K|=\pi/a$, and contribute equally to $\mathcal{S}$, whereas the local model retains only the first. The crossing wavenumbers, $k_ca/\pi=0.87$ and $0.61$, correspond to $\varepsilon_s=k_ca/(2\pi)=0.44$ and $0.31$ in equation~\eqref{eq:homogenization_scale_separation}. This is the regime in which local resonance and Bragg scattering interact, and in which the discrete formulation is required. {\color{red}The finite element model, which shares no part of the solution with the lattice sums, reproduces both edges to $10^{-10}$ or better (Section~\ref{subsubsec:num_L1_5}).}

\paragraph{Dependence on the lattice spacing.} Figure~\ref{fig:lattice_design}(a) shows the width of the complete stop band as a function of the spacing $a$, at fixed loading $n_0m_{\rho,r}=m_{\rho,r}/a=0.01$~m, computed from the lattice sums~\eqref{eq:bloch_condition}. Away from the Bragg regime, the lower edge is the local lower branch at the zone boundary, depressed twice as much as in the local model because the components $K=0$ and $K=-2\pi/a$ contribute equally, and the upper edge is $\omega_\Gamma$. This gives
\begin{equation}\label{eq:gap_asymptotic}
    \frac{\omega_\Gamma-\omega_X}{\omega_r}\simeq\frac{\mu(\pi/a)\,\omega_r^2}{\omega_0^2(\pi/a)-\omega_r^2}-\frac{\Pi_L}{2},
\end{equation}
with $\mu$ and $\omega_0$ the mass ratio and bare branch of Section~\ref{subsubsec:fluid_dilution}, which reproduces the exact widths within $2\%$ for $ak_c/\pi\lesssim0.75$. When bending dominates the host at the lattice scale, $\omega_0^2(\pi/a)\simeq C_\rho(\pi/a)^4/(m_\rho+M_{f,\rho})$ and the host inertia cancels, so that
\begin{equation}\label{eq:gap_a4}
    \frac{\omega_\Gamma-\omega_X}{\omega_r}\simeq\Big(\frac{1}{\pi^4}-\frac{1}{1440}\Big)\frac{n_0\,m_{\rho,r}\,\omega_r^2\,a^4}{C_\rho}.
\end{equation}
The complete stop band of a floating lattice therefore vanishes as $a^4$ in the homogenisation limit, $a\to0$ {\color{red}(finite element slope $4.01$ for $a\le0.5~\si{\meter}$, Section~\ref{subsubsec:num_L1_6})}. This is the opposite of a metaplate in vacuo, whose local-resonance gap tends to the finite value $\omega_r(\sqrt{1+\mu}-1)$, i.e. $0.41\,\omega_r$ for these resonators (dotted line in Figure~\ref{fig:lattice_design}(a)). The role of the fluid is to remove the long-wave part of the gap (Section~\ref{subsubsec:fluid_dilution}); what remains is a discreteness effect, which grows only when the zone boundary approaches the crossing wavenumber, $ak_c/\pi\to1$, i.e. in the regime $\varepsilon_s=O(1)$ where the local model does not apply. At $a=1$~m, the width $0.034\,\omega_r$ ($T_\rho=0$) is $1.8$ times the estimate obtained by truncating the local branches at $qa=\pi$, and $12$ times smaller than the in-vacuo value.

\begin{figure*}[pos=htbp]
    \centering
    \includegraphics[width=\linewidth]{figures/lattice_design.png}
    \caption{Hydroelastic limits of discrete resonators, parameters of Table~\ref{tab:dispersion_test_parameters}. (a) Complete stop band of a floating line array as a function of the spacing, at fixed loading $n_0m_{\rho,r}=0.01$~m. Symbols: exact lattice sums, equation~\eqref{eq:bloch_condition}; lines: equation~\eqref{eq:gap_asymptotic}; dotted: homogenised gap of the same resonators on a metaplate in vacuo. (b) Optimal damping ratio of a single resonator in the vertical section, equation~\eqref{eq:zeta_opt}, with $m_{\rho,r}=0.01~\si{\meter\squared}$; dotted: bound~\eqref{eq:zeta_bound}; shaded: ocean-wave band.}
    \label{fig:lattice_design}
\end{figure*}

\begin{figure*}[pos=htbp]
    \centering
    \includegraphics[width=0.85\linewidth]{figures/discrete_vs_homogenised_bands.png}
    \caption{Bloch bands of a line array of identical resonators (spacing $a=1$~m, parameters of Table~\ref{tab:dispersion_test_parameters}) for $T_\rho=0$ (left) and $T_\rho=9.8~\si{\meter\cubed\per\second\squared}$ (right). Dots: exact discrete lattice, equation~\eqref{eq:bloch_condition}. Dashed: local branches, equation~\eqref{eq:omega_k_x_closed_form}. Shaded: complete stop band $[\omega_X,\omega_\Gamma]$ of the discrete lattice. Dotted: $\omega_r$. The local model has no complete stop band.}
    \label{fig:lat_bands}
\end{figure*}

In configuration~B, the structure has finite length $L$ and the divergence of the added mass is cut off at the lowest wavenumber that the structure supports, $k\sim\pi/L$. To leading order, the attainable mass ratio \eqref{eq:lat_mu} is then $\mu_B\sim n_0m_{\rho,r}\,\pi\tanh(\pi H/L)/L$, and, by equation~\eqref{eq:anticrossing}, the hybridisation strength scales as $\sqrt{\mu_B}$. Unlike in the in-vacuo metaplate, it therefore depends on the size of the structure and decreases as $L$ increases.
%% [2026-10-28 numerical feedback, J2025-LRMM Sec. 6]
% TODO (verify the configuration-B estimate against the finite-element model): done, Sec. 6 L2-5.
{\color{red}The finite element computations of Section~\ref{subsubsec:num_L2_5} confirm the exponents of this estimate: the splitting decreases as $L^{-0.54}$ in deep water, and in shallow water its slope steepens towards $-1$ as $L/(\pi H)$ grows. The prefactor is less accurate. For five or more resonators, $\sqrt{\mu_B}$ underestimates the splitting by a factor of $1.7$--$1.8$ in deep and intermediate water, and the local mass ratio~\eqref{eq:lat_mu} at the modal wavenumber $\beta_1/L$, $\beta_1=4.730$, by $1.4$--$1.5$. The mass ratio formed with the wet modal mass of the bare mode gives the splitting to $1$--$6.5\%$.}

% \subsection{Line arrays: an exact multimode dispersion relation}\label{subsec:lat_line_arrays}

% In the vertical section, the residue form~\eqref{eq:lat_green_1d} and the identity
% $\sum_{n\in\mathbb{Z}}e^{ik_m|n|a-iqna}=i\sin(k_ma)/(\cos qa-\cos k_ma)$ turn the Bloch condition~\eqref{eq:lat_bloch} into
% \begin{equation}\label{eq:lat_1d_exact}
%     \frac{1}{\alpha(\omega)}=-\sum_m\frac{\sin(k_ma)}{\mathcal{D}_0'(k_m)\,\bigl(\cos qa-\cos k_ma\bigr)} ,
% \end{equation}
% where $m$ runs over the propagating root, the flexural roots and all evanescent fluid roots. Retaining only $k_0$ recovers the classical transfer-matrix law
% \begin{equation}\label{eq:lat_1d_single}
%     \cos qa=\cos k_0a-\frac{\alpha(\omega)}{\mathcal{D}_0'(k_0)}\sin k_0a ,
% \end{equation}
% which reduces to the mass-loaded string for $\mathcal{D}_0=T_\rho k^2-\omega^2m_\rho$. Equation~\eqref{eq:lat_1d_exact} quantifies when equation~\eqref{eq:lat_1d_single} fails. The neglected terms are $O(e^{-a\operatorname{Im}k_m})$ and become significant when the spacing is comparable with the decay length of the flexural roots or with $H/\pi$. At fixed real $\omega$, equation~\eqref{eq:lat_1d_exact} truncated to $M$ roots is a rational equation of degree $M$ in $X=\cos qa$. Its roots give all propagating and evanescent Bloch modes of the array, which is the basis for mode matching at the ends of finite arrays.

% For a 3D rectangular lattice (spacing $a$ along $x$ and $b$ along $y$), applying the same summation in $x$ after expanding the $y$-periodic Green function in diffraction orders $\beta_\nu=q_y+2\pi\nu/b$ gives
% \begin{equation}\label{eq:lat_2d_rect}
%     \frac{1}{\alpha(\omega)}=-\frac{1}{b}\sum_{\nu\in\mathbb{Z}}\sum_m
%     \frac{k_m\,\sin(\kappa_{\nu m}a)}{\kappa_{\nu m}\,\mathcal{D}_0'(k_m)\,\bigl(\cos q_xa-\cos\kappa_{\nu m}a\bigr)},
%     \qquad \kappa_{\nu m}=\sqrt{k_m^2-\beta_\nu^2},\quad\operatorname{Im}\kappa_{\nu m}\ge0 .
% \end{equation}
% For $|\nu|\to\infty$, $\kappa_{\nu m}\simeq i|\beta_\nu|$ and the summand tends to a $q_x$-independent asymptote that decays only algebraically in $\nu$. Subtracting this asymptote, which reproduces the self-term $G(0)$ and can be summed in closed form, leaves terms that decay like $e^{-|\beta_\nu|a}$.

\subsection{Finite array of resonators on infinite structure}\label{subsec:finite_array}

We now consider $N$ resonators at arbitrary positions $\xvec_i$, with possibly distinct attachment strengths $\alpha_i$, excited by an incident field $\widehat{\eta}_f$. Evaluating equation~\eqref{eq:green_solution} at the attachment points gives the Foldy system \citep{foldy1945multiple,martin2006multiple}
\begin{equation}\label{eq:lat_foldy}
    \bigl[\mathbf{I}-\mathbf{G}\boldsymbol{\alpha}\bigr]\boldsymbol{\eta}_a=\boldsymbol{\eta}_a^{\mathrm{inc}},
    \qquad
    G_{ij}=G(\xvec_i-\xvec_j),
    \qquad
    \boldsymbol{\alpha}=\operatorname{diag}(\alpha_1,\ldots,\alpha_N),
\end{equation}
where $\boldsymbol{\eta}_a$ and $\boldsymbol{\eta}_a^{\mathrm{inc}}$ collect the total and incident deflections at the attachments, and $G_{ii}=G(0)$. Once $\boldsymbol{\eta}_a$ is known, the full field development follows from equation~\eqref{eq:green_solution}. Equation~\eqref{eq:lat_foldy} is exact within the point-attachment model and has four properties that are used below. First, it retains all the interactions transmitted through the host, including the evanescent near field of the flexural and fluid roots. Second, $\mathbf{G}$ depends only on the host structure and on the attachment positions, whereas the resonators enter only through the diagonal matrix $\boldsymbol{\alpha}$. Therefore, graded resonator populations, detuned resonators or parametric studies only modify $\boldsymbol{\alpha}$. Third, reciprocity of the host implies $\mathbf{G}=\mathbf{G}^{\mathsf{T}}$, i.e. $\mathbf{G}$ is complex symmetric but not Hermitian. Fourth, the free modes of the array are the non-trivial solutions of the homogeneous system, $\det[\mathbf{I}-\mathbf{G}\boldsymbol{\alpha}]=0$. For a finite array these occur at complex frequencies, since energy leaks through the ends of the array, and for an infinite periodic array with Bloch phases, equation~\eqref{eq:lat_foldy} reduces to the Bloch condition \eqref{eq:bloch_condition}.

\subsubsection{Energy flux of finite array of resonators}
Here we derive the energy identities for three different configurations of a finite array of resonators. The energy identities for all the configurations below follow from a single power balance. Let $\mathbf{F}=\boldsymbol{\alpha}\boldsymbol{\eta}_a$ be the vector of resonator forces, so that $\boldsymbol{\eta}_a=\boldsymbol{\eta}_a^{\mathrm{inc}}+\mathbf{G}\mathbf{F}$. As in Section~\ref{subsubsec:energy_flux_single}, the power absorbed by the resonators is the work done by the structure on them, $P_{\mathrm{abs}}=\tfrac12\omega\operatorname{Im}(\mathbf{F}^{\mathsf{T}}\overline{\boldsymbol{\eta}_a})$. Using the symmetry of $\mathbf{G}$, this gives
\begin{equation}\label{eq:lat_power_balance}
    P_{\mathrm{abs}}
    =\tfrac12\omega\sum_j\gamma_{d,j}|F_j|^2
    =\tfrac12\omega\Bigl[\operatorname{Im}\bigl(\boldsymbol{\eta}_a^{\mathrm{inc}\,\mathsf{H}}\mathbf{F}\bigr)-\mathbf{F}^{\mathsf{H}}\mathbf{B}\,\mathbf{F}\Bigr],
    \qquad
    \mathbf{B}=\operatorname{Im}\mathbf{G},
\end{equation}
where the first equality is the power dissipated in the dashpots, with $\gamma_{d,j}=-\operatorname{Im}(1/\alpha_j)$. The last term, $\tfrac12\omega\mathbf{F}^{\mathsf{H}}\mathbf{B}\mathbf{F}$, is the power radiated by the forces $\mathbf{F}$, generalising equation~\eqref{eq:point_power}, and the first term is the power extracted from the incident wave. The real symmetric matrix $\mathbf{B}$ is the radiation matrix of the array. Maximising equation~\eqref{eq:lat_power_balance} over $\mathbf{F}$ gives the optimal forces $\mathbf{F}_{\mathrm{opt}}=\tfrac{\mathrm{i}}{2}\mathbf{B}^{+}\boldsymbol{\eta}_a^{\mathrm{inc}}$ and the maximum absorbed power
\begin{equation}\label{eq:lat_pmax}
    P_{\mathrm{abs}}\le P_{\max}=\frac{\omega}{8}\,\boldsymbol{\eta}_a^{\mathrm{inc}\,\mathsf{H}}\,\mathbf{B}^{+}\,\boldsymbol{\eta}_a^{\mathrm{inc}},
\end{equation}
where $\mathbf{B}^{+}$ is the pseudo-inverse of $\mathbf{B}$. This is the hydroelastic counterpart of the classical result for arrays of wave-energy converters \citep{falnes1980radiation}, with the radiation matrix given directly by the imaginary part of the host Green function. For $N=1$, $\mathbf{B}=\gamma_r$ and equation~\eqref{eq:lat_pmax} recovers the bound \eqref{eq:absorption_bound}. The bound is attained only if the forces $\mathbf{F}_{\mathrm{opt}}$ can be realised, which for passive, identical resonators is generally not the case.

\paragraph{Case 1. Line array in the vertical section.}
In the vertical section, the structure is one-dimensional and the Green function follows from the same residue calculus as equation~\eqref{eq:green_function},
\begin{equation}\label{eq:lat_green_1d}
    G(x)=\mathrm{i}\sum_m\frac{\mathrm{e}^{\mathrm{i}k_m|x|}}{\mathcal{D}_0'(k_m)},
    \qquad
    \operatorname{Im}G(x)=\frac{\cos(k_0x)}{\mathcal{D}_0'(k_0)}\quad\text{(non-dissipative)},
\end{equation}
where $m$ runs over the same roots as in equation~\eqref{eq:green_function}. A unit point force now radiates two waves of amplitude $1/\mathcal{D}_0'(k_0)$, one in each direction, and their total flux, given by equation~\eqref{eq:energy_flux}, equals $\tfrac12\omega\operatorname{Im}G(0)$, in agreement with equation~\eqref{eq:point_power}. For an incident wave $A\mathrm{e}^{\mathrm{i}k_0x}$, the propagating term of equation~\eqref{eq:lat_green_1d} gives, as $x\to\mp\infty$, the reflection and transmission coefficients
\begin{equation}\label{eq:lat_RT}
    \mathcal{R}=\frac{\mathrm{i}}{\mathcal{D}_0'(k_0)A}\sum_j\alpha_j\widehat{\eta}_j\,\mathrm{e}^{\mathrm{i}k_0x_j},
    \qquad
    \mathcal{T}=1+\frac{\mathrm{i}}{\mathcal{D}_0'(k_0)A}\sum_j\alpha_j\widehat{\eta}_j\,\mathrm{e}^{-\mathrm{i}k_0x_j}.
\end{equation}
In this case $B_{ij}=[\cos k_0x_i\cos k_0x_j+\sin k_0x_i\sin k_0x_j]/\mathcal{D}_0'(k_0)$ has rank two, one for each radiation direction, and the radiated power in equation~\eqref{eq:lat_power_balance} is exactly the flux of the scattered parts of $\mathcal{R}$ and $\mathcal{T}$. Dividing equation~\eqref{eq:lat_power_balance} by the incident flux \eqref{eq:energy_flux} then gives the exact identity
\begin{equation}\label{eq:lat_energy_1d}
    |\mathcal{R}|^2+|\mathcal{T}|^2+\frac{2}{\mathcal{D}_0'(k_0)|A|^2}\sum_j\operatorname{Im}\alpha_j\,|\widehat{\eta}_j|^2=1,
\end{equation}
where the last term is the fraction absorbed by the dashpots, since $\gamma_{d,j}|F_j|^2=\operatorname{Im}\alpha_j|\widehat{\eta}_j|^2$. The evanescent and complex roots contribute to $\mathbf{G}$, and thus to the response, but carry no flux to infinity, so equation~\eqref{eq:lat_energy_1d} is a two-channel balance regardless of the number of roots retained. For a single resonator, solving equation~\eqref{eq:lat_foldy} as in equation~\eqref{eq:single_force}, with $\gamma_r=1/\mathcal{D}_0'(k_0)$, gives the absorbed fraction
\begin{equation}\label{eq:lat_absorbed_1d}
    \mathcal{A}_1=\frac{2\gamma_r\gamma_d}{x^2+(\gamma_r+\gamma_d)^2}\le\frac12,
\end{equation}
with equality at critical coupling. This is the classical limit of one half for a symmetric absorber in two dimensions \citep{evans1976theory}. For $N\ge2$ resonators at distinct positions, equation~\eqref{eq:lat_pmax} gives $P_{\max}$ equal to the incident flux, i.e. complete absorption is possible in principle, since two forces can control both channels.

For identical resonators with spacing $a$, the transmission through the array is related to its Bloch spectrum. If only the propagating root is retained in the lattice sum of equation~\eqref{eq:bloch_condition}, the sum can be evaluated in closed form, giving the single-mode dispersion relation
\begin{equation}\label{eq:lat_1d_single}
    \cos qa=\cos k_0a-\frac{\widehat{\alpha}}{\mathcal{D}_0'(k_0)}\sin k_0a,
    \qquad
    \frac{1}{\widehat{\alpha}}=\frac{1}{\alpha}-\operatorname{Re}G(0),
\end{equation}
which reduces to the classical mass-loaded string for $\mathcal{D}_0=Tk^2-m\omega^2$. Here, the evanescent part of the self-interaction has been kept through the renormalised strength $\widehat{\alpha}$, as in equation~\eqref{eq:renorm_single}, therefore equation~\eqref{eq:lat_1d_single} is exact for $N=1$. The evanescent interactions between neighbouring resonators are of order $\mathrm{e}^{-a\operatorname{Im}k_m}$. These interactions are neglected here, thus, equation~\eqref{eq:lat_1d_single} is accurate when the spacing is large compared with the decay length of the flexural roots and with $H/\pi$. Within this approximation, each resonator acts as a non-dissipative scatterer with $1/|\mathcal{T}_1|^2-1=(\widehat{\alpha}/\mathcal{D}_0'(k_0))^2$, and the standard theory of periodic stacks gives, for $N$ identical lossless resonators,
\begin{equation}\label{eq:lat_fabry_perot}
    \frac{1}{|\mathcal{T}_N|^2}=1+\Big(\frac{\widehat{\alpha}}{\mathcal{D}_0'(k_0)}\Big)^2\frac{\sin^2(Nqa)}{\sin^2(qa)} .
\end{equation}
Equation~\eqref{eq:lat_fabry_perot} makes the link between the Bloch spectrum and finite-array transmission explicit. In a pass band, $q$ is real, $|\mathcal{T}_N|\le1$ oscillates with $N$, and $|\mathcal{T}_N|=1$ at the Fabry--P\'erot resonances $Nqa=p\pi$, $p=1,\ldots,N-1$. In a stop band, $q=\mathrm{i}\varkappa$ or $q=\pi/a+\mathrm{i}\varkappa$, with $\varkappa>0$, and for $N\varkappa a\gg1$
\begin{equation}\label{eq:lat_decay}
    |\mathcal{T}_N|\simeq\frac{2\sinh(\varkappa a)\,\mathcal{D}_0'(k_0)}{|\widehat{\alpha}|}\,\mathrm{e}^{-N\varkappa a},
\end{equation}
i.e. the transmission decays exponentially with the number of resonators, at the rate given by the imaginary part of the Bloch wavenumber, and with a prefactor set by the single-resonator reflectivity. This justifies the use of complex Bloch wavenumbers at real frequency, Section~\ref{subsec:infinite_array}, to characterise forced finite arrays. {\color{red}In finite depth, long arrays have a second, weakly attenuated Bloch mode, which continues the slowest evanescent fluid mode. It dominates the transmission wherever the propagating attenuation exceeds its own, so that the attenuation per unit length of a long floating array is bounded by approximately $\kappa_1\simeq\pi/(2H)$ (Section~\ref{subsubsec:num_L1_4}).}

The corresponding identities for a grating stack, i.e. rows periodic in $y$ under oblique incidence, and for a finite cluster in an unbounded plate, including the optical theorem and the maximum absorption cross-section of the cluster, are given in Appendix~\ref{app:lattice_supp}.

\subsection{Finite structures in open water}\label{sec:finite_domain}

For a structure of finite extent floating in open water, the host is no longer translation invariant and the Green function must be replaced by the transfer compliance of the \emph{bare} structure, $G_B(\xvec_i,\xvec_j;\omega)$, i.e. the deflection at $\xvec_i$ due to a unit point force at $\xvec_j$, with the radiation condition imposed in the open water. Equations~\eqref{eq:lat_foldy}--\eqref{eq:lat_pmax} hold unchanged with $G_{ij}=G_B(\xvec_i,\xvec_j)$, since they only rely on linearity and reciprocity. Unlike in configuration~A, where the fluid enters the symbol \eqref{eq:host_symbol} through a frequency-independent added mass, $G_B$ depends on frequency through the open-water wavenumber, the radiation condition and the edge conditions \eqref{eq:structure_bcs}. This is the frequency-dependent hydrodynamic operator referred to in Section~\ref{sec:dispersion_relation}. Moreover, $\operatorname{Im}G_B$ now represents radiation of free-surface waves into the open water, so that a finite structure radiates even when its resonators are arranged periodically.

The matrix $\mathbf{G}_B$ can be obtained from the finite-element model of Section~\ref{sec:numerical_formulation} with the resonator unknowns removed. Each column requires the solution for one right-hand side, a unit load in the structural equation at an attachment point, and all columns share one factorisation per frequency. The connection with the monolithic model is exact at the discrete level. Let $\mathbf{K}_h(\omega)$ be the discrete operator of the bare fluid--structure problem and $\mathbf{P}=[\mathbf{p}_1,\ldots,\mathbf{p}_N]$, where $\mathbf{p}_i$ evaluates the structural trace at $\xvec_i$. Eliminating the resonator unknowns $q_i$ from the monolithic system adds the rank-$N$ term $-\mathbf{P}\boldsymbol{\alpha}\mathbf{P}^{\mathsf{T}}$ to $\mathbf{K}_h$. By the Sherman--Morrison--Woodbury identity, the attachment deflections of the monolithic solution, $\boldsymbol{\eta}_a=\mathbf{P}^{\mathsf{T}}(\mathbf{K}_h-\mathbf{P}\boldsymbol{\alpha}\mathbf{P}^{\mathsf{T}})^{-1}\mathbf{f}_h$, then satisfy
\begin{equation}\label{eq:lat_woodbury}
    \bigl[\mathbf{I}-\mathbf{G}_{B,h}\boldsymbol{\alpha}\bigr]\boldsymbol{\eta}_a=\boldsymbol{\eta}_a^{\mathrm{inc}},
    \qquad
    \mathbf{G}_{B,h}=\mathbf{P}^{\mathsf{T}}\mathbf{K}_h^{-1}\mathbf{P},
    \qquad
    \boldsymbol{\eta}_a^{\mathrm{inc}}=\mathbf{P}^{\mathsf{T}}\mathbf{K}_h^{-1}\mathbf{f}_h,
\end{equation}
which is equation~\eqref{eq:lat_foldy} with the discrete transfer compliance $\mathbf{G}_{B,h}$, \emph{to round-off}. Hence, once $\mathbf{G}_{B,h}$ and $\boldsymbol{\eta}_a^{\mathrm{inc}}$ are computed, any choice of resonator parameters, including grading and detuning, can be evaluated at the cost of an $N\times N$ solve, and the eigenvalue problem $\det[\mathbf{I}-\mathbf{G}_{B,h}\boldsymbol{\alpha}]=0$ gives the coupled modes of the structure with resonators. Reciprocity, $\mathbf{G}_B=\mathbf{G}_B^{\mathsf{T}}$, holds for the continuous problem. It also holds for $\mathbf{G}_{B,h}$, to round-off and for any mesh: $\mathbf{K}_h$ is not symmetric, but changing the sign of its potential rows makes it complex-symmetric without changing the solution (Remark~\ref{rem:discrete_reciprocity}). This is no longer true when free-surface damping zones are used, so that the residual $\|\mathbf{G}_{B,h}-\mathbf{G}_{B,h}^{\mathsf{T}}\|$ diagnoses the absorbing-layer treatment rather than the discretisation. {\color{red}With the Sommerfeld condition, reciprocity holds to $1.7\times10^{-9}$ or better and the Woodbury identity~\eqref{eq:lat_woodbury} to $2.6\times10^{-10}$ or better (Section~\ref{subsubsec:num_L2_1}).}

The single-resonator results of Section~\ref{subsubsec:energy_flux_single} carry over with $G(0)$ replaced by $G_B(\xvec_1,\xvec_1)$, so that $\gamma_r=\operatorname{Im}G_B(\xvec_1,\xvec_1)$ and the renormalised resonance \eqref{eq:renorm_single} now depends on the attachment position. In the vertical section, let $a^\pm$ be the amplitudes of the waves radiated towards $x\to\pm\infty$ by a unit force at $\xvec_1$. The radiated power, and hence $\gamma_r$, is proportional to $|a^+|^2+|a^-|^2$. By reciprocity, the deflection induced at $\xvec_1$ by a wave incident from $x\to-\infty$ is proportional to $a^-$, with the same constant, a Haskind-type relation. Inserting both into equation~\eqref{eq:lat_pmax} with $N=1$ bounds the absorbed fraction by
\begin{equation}\label{eq:lat_absorbed_B}
    \mathcal{A}\le\frac{|a^-|^2}{|a^-|^2+|a^+|^2},
\end{equation}
which equals $\tfrac12$ for a resonator at the centre of a symmetric structure, in agreement with equation~\eqref{eq:lat_absorbed_1d} and with the classical result for two-dimensional absorbers \citep{evans1976theory}. Equation~\eqref{eq:lat_absorbed_B} identifies asymmetric radiation, e.g. through an off-centre attachment or an asymmetric structure, as the route by which a single resonator can absorb more than one half of the incident power. {\color{red}On a free plate, the finite element absorbed fraction at critical coupling equals the bound~\eqref{eq:lat_absorbed_B} to $1.8\times10^{-4}$, with values between $0.088$ and $0.912$ along the plate (Section~\ref{subsubsec:num_L2_4}).} With $N\ge2$ resonators, the radiation matrix has rank two and, as for the line array of Section~\ref{subsec:finite_array} (Case~1), complete absorption becomes possible in principle.